\documentclass[../main.tex]{subfiles}
\begin{document}
\begin{center}
    \Large{\textbf{DANH MỤC KÝ HIỆU}}
\end{center}
\phantomsection\addcontentsline{toc}{chapter}{DANH MỤC KÝ HIỆU}
\vspace{0.3cm}

\begin{center}
\renewcommand{\arraystretch}{1.25}
\begin{tabular}{@{}p{3cm}p{10.5cm}@{}}
\toprule
\textbf{Ký hiệu} & \textbf{Ý nghĩa} \\
\midrule
$F_s$, $T_s$ & Tần số lấy mẫu (Hz) và chu kỳ lấy mẫu (s) \\
$N$ & Số mẫu của tín hiệu \\
$f_N=F_s/2$ & Tần số Nyquist \\
$s(n)$ & Tín hiệu hữu ích (tín hiệu chuẩn trong mô phỏng) \\
$v(n)$ & Nhiễu cộng, gồm rung động cơ $v_m(n)$ và nhiễu ngẫu nhiên $v_r(n)$ \\
$d(n)$ & Tín hiệu đo (tín hiệu mong muốn của bộ lọc), $d(n)=s(n)+v(n)$ \\
$x(n)$, $\vect{x}(n)$ & Tín hiệu tham chiếu và vectơ hồi quy $M$ mẫu \\
$y(n)$ & Đầu ra bộ lọc FIR thích nghi \\
$e(n)$ & Sai số $e(n)=d(n)-y(n)$ \\
$\vect{w}(n)$ & Vectơ trọng số bộ lọc ($M\times 1$) \\
$M$ & Bậc (số hệ số) của bộ lọc \\
$\mu$ & Bước học của LMS/NLMS \\
$\varepsilon$ & Hằng số chuẩn hoá của NLMS \\
$\lambda$ & Hệ số quên của RLS \\
$\delta$ & Hằng số khởi tạo $\vect{P}(0)=\delta^{-1}\vect{I}$ của RLS \\
$\vect{P}(n)$ & Ước lượng nghịch đảo ma trận tự tương quan ($M\times M$) \\
$\vect{k}(n)$ & Vectơ độ lợi của RLS \\
$\Delta$ & Độ trễ giải tương quan của cấu trúc ALE (mẫu) \\
$\vect{R}_{xx}$, $\lambda_{\max}$ & Ma trận tự tương quan đầu vào và trị riêng lớn nhất của nó \\
\bottomrule
\end{tabular}
\end{center}

\end{document}
